\documentclass[a4paper,10pt]{article}
\usepackage[utf8]{inputenc}
\RequirePackage{amsthm,amsmath,amssymb,amscd,xcolor}
\def\D{^{C}D_t^\alpha}
\def\W{\mathbb{W}}
\def\H{\mathbb{H}}
\def\K{\mathbb{K}}
\def\E{\mathbb{E}}
\def\cL{\mathcal{L}}
\def\R{\mathbb{R}}
\def\ud{\text{d}}
\newcommand{\Norm}[1]{\left|\left|  #1   \right|\right|}
\begin{document}

Here are comments for {\it d1171e49} on the file -- {\it thesis.tex}:

\begin{enumerate}
  \item The line after (1.2): and $\ln_+(z) \textcolor{magenta}{:= }...$. 
  \item The paragraph after the centered equation for $I(t,x)$ on p. 2:
    \begin{enumerate}
      \item ``\textcolor{magenta}{consists of} a triple";
      \item Fox-H function should be italic, i.e., {\it  Fox-H functions};
      \item ``where $\mathcal{F}f(\xi) ...$ \textcolor{magenta}{refers to the Fourier transform.}";
      \item ``$a>0,\: b>0, \: \textcolor{magenta}{z\in \mathbb{C}},$";
      \item Remove the hyphen between {\it Gamma-function}.
    \end{enumerate}
  \item In the expression for $\mathcal{T}_0$ on p. 5, for the case $\beta\in (0,1]$ (resp. $\beta\in (1,2)$), replace  $
    \left\lceil \beta \right\rceil -1 $ by $0$ (resp. $1$). Carry out the same simplification in the proof for these
    cases. Do the simplification whenever possible and Check if you have similar things throughout the draft.
  \item Statement of Lemma 2.2, the last sentence should better be updated as ``Moreover, when $\beta\in (1,2)$, ".
  \item (2.9): Remove $t\in [0,T]$ and write $\displaystyle \sup_{(t,x)\in[0,T]\times\R^d}$ instead.
  \item The centered equation before (2.11): $(u)^{\rho-1}$ $\to$  $u^{\rho-1}$. The sentence afterwards, [CHN19] should be
    (2.4). 
  \item The sentence after (2.11): ``\textcolor{magenta}{if} we now let ". Add a comma after the next centered equation.
  \item Line after (2.13): ``It is important to note" $\to$ ``Note that".
  \item Proposition 2.5:
    \begin{enumerate}
      \item In the statement: [CH21, Proposition\textcolor{magenta}{s} 2.1 and 2.2].
      \item In the statement, no need to write again $Y(t,x) = ...$. We already know what is  $Y(t,x)$.
      \item In the statement of case 1, add ``, then " before ``there". The same for case 2.
    \end{enumerate}
  \item Proposition 2.6:
    \begin{enumerate}
      \item In the statement, Remove ``let $a>0$ and" and put it afterwards: ``Then for all $a>0$, $x,z\in \R$ ,....".
      \item[Note:] I prefer to use ``all" instead of ``any". It is less ambiguous. I learnt this from a math writing
        class at EPFL.
      \item Rewrite the last two lines on p. 7 as:
        \begin{center}
          and $q\in (0,\rho]$ (resp. $q\in (0,\rho)$) under Case 1 (resp. 2) of Proposition 2.5.
        \end{center}
        Do the same for the statement of Proposition 2.7.
    \end{enumerate}
  \item Statement of Proposition 2.7: ``Moreover, if we consider the case when".
  \item When using the capitalized word, no need to use ``the" in front: Remove the second ``the" at the last line of
    the proof of Theorem 1.2.
\end{enumerate}

\end{document}
